\documentclass{amsart}
\usepackage{amsmath, amssymb, amsfonts, amsthm}
\usepackage{datetime}
\usepackage{graphicx}
\usepackage[
colorlinks=true,
linkcolor=blue,
citecolor=magenta,
urlcolor=red
]{hyperref}
%\usepackage{kpfonts, euler}
\usepackage{url}
\usepackage{newtxtext, mathpazo}
\numberwithin{equation}{section}
\theoremstyle{plain}
\newtheorem{theorem}{Theorem}[section]
\newtheorem{lemma}[theorem]{Lemma}
\newtheorem{proposition}[theorem]{Proposition}
\newtheorem{corollary}[theorem]{Corollary}
\theoremstyle{definition}
\newtheorem{definition}[theorem]{Definition}
\newtheorem{example}[theorem]{Example}
\newtheorem{remark}[theorem]{Remark}
\newcommand{\N}{\mathbb{N}}
\newcommand{\Z}{\mathbb{Z}}
\newcommand{\Q}{\mathbb{Q}}
\newcommand{\R}{\mathbb{R}}
\newcommand{\F}{\mathbb{F}}
\newcommand{\powerset}{\mathcal{P}}


\title{M\"obius Inversion through \\ Partially Ordered Sets}
\author{Sayan Kar}
\date{Last updated: \today}


\begin{document}

\maketitle
\tableofcontents

\section{M\"obius Function}

\begin{definition}
A \emph{partially ordered set} (or \emph{poset}) is a pair $(P,\preceq)$ where $P$ is a set and $\preceq$ is a binary relation on $P$ satisfying the following properties for all $x,y,z\in P$:
\begin{enumerate}
    \item \textbf{Reflexivity:} $x \preceq x$.
    \item \textbf{Antisymmetry:} if $x \preceq y$ and $y \preceq x$, then $x=y$.
    \item \textbf{Transitivity:} if $x \preceq y$ and $y \preceq z$, then $x \preceq z$.
\end{enumerate}
\end{definition}

\begin{lemma}
\label{topo_ordering}
Suppose that $|P|=n$. Then we can label the elements of $P=\{p_1,p_2,\ldots,p_n\}$ so that
\[
p_i \preceq p_j \implies i \le j.\footnote{This type of ordering of a poset is called a \emph{topological ordering}.}
\]
\end{lemma}

\begin{proof}
We prove this by induction on $n$. The base case $n=1$ is trivial. Choose a maximal element of $P$ and label it $p_n$. Let $P' = P \setminus \{p_n\}.$ By the induction hypothesis, we can label $P'=\{p_1,p_2,\ldots,p_{n-1}\}$ so that
\[
p_i \preceq p_j \implies i \le j.
\]
Since $p_n$ is maximal, this labeling together with $p_n$ gives the required labeling of all of $P$.
\end{proof}

Let $(P,\preceq)$ be a finite partially ordered set. We define the function
\[
\zeta : P^2 \to \{0,1\}
\]
by
\[
\zeta(x,y)=
\begin{cases}
1 & \text{if } x \preceq y,\\
0 & \text{otherwise}.
\end{cases}
\]

\begin{definition}
Let $P$ be a finite poset. Consider the matrix
\[
A_{\zeta} = [\zeta(x,y)]_{x,y\in P}.
\]
By Lemma~\ref{topo_ordering}, there is a topological ordering of $P$, the matrix $A_{\zeta}$ can an upper triangular matrix with all diagonal entries equal to $1$.
Hence $A_{\zeta}$ is invertible. Its inverse is denoted
\[
A_{\mu} = [\mu(x,y)]_{x,y\in P}.
\]
The function $\mu : P^2 \to \mathbb{Z}$ is called the \emph{M\"obius function} of the poset $P$.
\end{definition}

Since $A_{\mu}A_{\zeta}=I$, we obtain
\[
\sum_{z\in P}\mu(x,z)\zeta(z,y)
=
\sum_{x\preceq z\preceq y}\mu(x,z)
=
\delta(x,y).
\]

\begin{theorem}

Consider the two following cases:


\begin{enumerate}
\item[(a)] Let $P$ be the set of subsets of a finite set $X$, ordered by inclusion $\subseteq$. Then
\[
\mu(A,B)=
\begin{cases}
(-1)^{|B\setminus A|} & \text{if } A\subseteq B,\\
0 & \text{otherwise}.
\end{cases}
\]

\item[(b)] Let $P=[n]$ with the order defined by $a \preceq b$ if $a$ divides $b$. Then
\[
\mu(a,b)=
\begin{cases}
(-1)^r & \text{if } b/a \text{ is the product of } r \text{ distinct primes},\\
0 & \text{otherwise}.
\end{cases}
\]
\end{enumerate}
\end{theorem}

\begin{proof}
We only need to verify the defining relation
\[
\sum_{x \preceq z \preceq y} \mu(x,z) =
\begin{cases}
1 & x=y,\\
0 & x\ne y.
\end{cases}
\]

\begin{enumerate}
\item[(a)] If $A\subseteq B$, then
\[
\sum_{A\subseteq C\subseteq B} (-1)^{|C\setminus A|}
= \sum_{k=0}^{|B\setminus A|} \binom{|B\setminus A|}{k} (-1)^k
= (1-1)^{|B\setminus A|}.
\]
This equals $0$ unless $A=B$, in which case it equals $1$.

\item[(b)] Suppose
\[
b/a = p_1^{k_1}p_2^{k_2}\cdots p_r^{k_r},
\]
where $p_1,p_2,\ldots,p_r$ are distinct primes and $k_1,k_2,\ldots,k_r \ge 1$.
Then any $c$ satisfying $a\mid c\mid b$ corresponds to choosing a subset
of the primes $\{p_1,\ldots,p_r\}$. Hence
\[
\sum_{a\mid c\mid b} \mu(c,b)
= \sum_{S\subseteq [r]} (-1)^{|S|}
=
\begin{cases}
1 & r=0,\\
0 & r\ge 1.
\end{cases}
\]
This verifies the required identity.
\end{enumerate}
\end{proof}

\section{M\"obius Inversion}
\begin{theorem}
Let $(P,\preceq)$ be a finite poset and let $f,g,h : P \to \mathbb{R}$ be functions such that
\begin{equation}
g(x)=\sum_{a \preceq x} f(a), \qquad 
h(x)=\sum_{b \succeq x} f(b).
\end{equation}
Then the function $f$ can be recovered from $g$ or $h$ via
\begin{equation}
f(x)=\sum_{a \preceq x}\mu(a,x)g(a), \qquad
f(x)=\sum_{b \succeq x}\mu(x,b)h(b).
\end{equation}
\end{theorem}

\begin{proof}
Since $P$ is finite, we may label its elements as $P=\{p_1,\dots,p_n\}$ in a way compatible with the partial order. Any function $f:P\to\mathbb{R}$ can then be identified with a column vector
\[
\mathbf{f} = (f(p_1),f(p_2),\ldots,f(p_n))^{T}.
\]
Similarly we associate vectors $\mathbf{g}$ and $\mathbf{h}$ to the functions $g$ and $h$.

Let $A_{\zeta}=[\zeta(x,y)]$ be the matrix corresponding to the zeta function of the poset. The relation
\[
g(x)=\sum_{a\preceq x}f(a)
\]
means that each coordinate of $\mathbf{g}$ is obtained by summing the entries of $\mathbf{f}$ over elements below it in the order. In matrix form this is
\[
\mathbf{g}=A_{\zeta}^{T}\mathbf{f}.
\]
Likewise, the relation
\[
h(x)=\sum_{b\succeq x}f(b)
\]
corresponds to the matrix equation
\[
\mathbf{h}=A_{\zeta}\mathbf{f}.
\]

Since $A_{\zeta}$ is invertible and its inverse is $A_{\mu}$, we obtain
\[
\mathbf{f}=A_{\zeta}^{-T}\mathbf{g}=A_{\mu}^{T}\mathbf{g}
\quad\text{and}\quad
\mathbf{f}=A_{\zeta}^{-1}\mathbf{h}=A_{\mu}\mathbf{h}.
\]
Writing these matrix identities co-ordinatewise yields the stated inversion formulas.
\end{proof}

\section{Applications of M\"obius Inversion}
\subsection{Inclusion--Exclusion Principle}

M\"obius inversion gives a particularly elegant proof of the classical
principle of inclusion--exclusion. Let $\{A_i\}_{i\in I}$ be a family of subsets of a finite set $X$.
For $J\subseteq I$, define
\[
f(J)=\left|\left(\bigcap_{i\in J}A_i\right)\cap
\left(\bigcap_{i\notin J}(X\setminus A_i)\right)\right|,
\]
that is, the number of elements that lie in exactly the sets indexed by $J$.
Let
\[
h(J)=\left|\bigcap_{i\in J}A_i\right|.
\]
Then
\[
h(J)=\sum_{K\supseteq J} f(K).
\]

Applying M\"obius inversion on the poset $(\mathcal{P}(I),\subseteq)$, we obtain
\[
f(J)=\sum_{K\supseteq J}\mu(J,K)h(K)
     =\sum_{K\supseteq J}(-1)^{|K|-|J|}h(K).
\]

Putting $J=\varnothing$ gives
\[
\left|\bigcap_{i\in I}(X\setminus A_i)\right|
=
\sum_{K\subseteq I}(-1)^{|K|}
\left|\bigcap_{j\in K}A_j\right|.
\]


\subsection{Totient Function}

M\"obius inversion gives another proof of the classical closed form
for Euler's totient function. We briefly note that the divisibility relation on the natural numbers
forms a partially ordered set. In this \emph{divisibility poset}, we write
$a \preceq b$ if $a$ divides $b$. For a function $f:\mathbb{N}\to\mathbb{R}$,
define
\[
g(n)=\sum_{d\mid n} f(d).
\]
M\"obius inversion on this poset gives
\[
f(n)=\sum_{d\mid n}\mu(d,n)g(d)
   =\sum_{\substack{d\mid n\\ n/d\ \text{square free}}}
     (-1)^{p(n/d)}\,g(d),
\]
where $p(m)$ denotes the number of distinct prime factors of $m$. For a natural number $n$, let $\varphi(n)$ denote the number of integers
$m\le n$ such that $\gcd(m,n)=1$.

\begin{lemma}
\[
n=\sum_{d\mid n}\varphi(d)=\sum_{d\mid n}\varphi(n/d).
\]
\end{lemma}

\begin{proof}
If $(m,n)=d$, then we can write $m=m_1d$ and $n=n_1d$ where
$\gcd(m_1,n_1)=1$. Hence the number of such $m$ is the number
of choices of $m_1$, which is $\varphi(n_1)=\varphi(n/d)$.
\end{proof}

Applying M\"obius inversion to this identity with $g(n)=n$
and $f(n)=\varphi(n)$ gives
\[
\varphi(n)=\sum_{d\mid n}(-1)^{p(n/d)}d
         =\sum_{d\mid n}(-1)^{p(d)}\frac{n}{d}.
\]
Factoring out $n$ we obtain
\[
\varphi(n)=n\sum_{d\mid n}\frac{(-1)^{p(d)}}{d}.
\]

If $n=p_1^{k_1}p_2^{k_2}\cdots p_r^{k_r}$ where
$p_1,p_2,\ldots,p_r$ are distinct primes and $k_1,k_2,\ldots,k_r\ge1$,
the multiplicative structure of the sum yields the familiar product form
\[
\varphi(n)=n\prod_{i=1}^{r}\left(1-\frac{1}{p_i}\right).
\]


\subsection{Colored Necklace Problem}

M\"obius inversion also appears naturally in many counting problems.
A classical example is the problem of counting colored necklaces.

Suppose we have an infinite supply of beads of $\lambda$ colours and
wish to form a necklace with $n$ beads. Any necklace can be represented
by a sequence $(a_1,\ldots,a_n)$, where we identify sequences that differ
by a cyclic permutation, since they represent the same necklace.

We say that a necklace of length $n$ is \emph{primitive} if it is not
obtained by repeating a shorter necklace. More precisely, if $d\mid n$,
a necklace is said to have period $d$ if it is obtained by repeating
$n/d$ times a primitive necklace of length $d$.

First note that the number of sequences of length $n$ is $\lambda^n$.
Every such sequence arises from a primitive necklace of some period
$d$ dividing $n$. If $M(d)$ denotes the number of primitive necklaces
of length $d$, then each such necklace gives rise to $d$ distinct
sequences depending on the choice of the starting position. Hence
\[
\lambda^n=\sum_{d\mid n} d\,M(d).
\]

Applying M\"obius inversion to this identity gives
\[
M(n)=\frac{1}{n}\sum_{d\mid n}\mu(d)\lambda^{n/d}.
\]

The total number of necklaces of length $n$ is therefore
\[
N_n=\sum_{d\mid n} M(d).
\]

Substituting the expression for $M(d)$, we obtain
\[
N_n=\sum_{d\mid n}\sum_{\ell\mid d}\frac{1}{d}\mu(\ell,d)\lambda^{\ell}.
\]
Writing $d=k\ell$ and rearranging the sums gives
\[
N_n=\sum_{\ell\mid n}\frac{\lambda^{\ell}}{\ell}
\sum_{k\mid n/\ell}\frac{\mu(1,k)}{k}.
\]

Using the identity
\[
\sum_{k\mid m}\frac{\mu(1,k)}{k}=\frac{\varphi(m)}{m},
\]
we finally obtain the well-known formula
\[
N_n=\frac{1}{n}\sum_{d\mid n}\varphi(n/d)\lambda^{d}.
\]

Thus the number of necklaces of length $n$ formed from $\lambda$
colours is
\[
N_n=\frac{1}{n}\sum_{d\mid n}\varphi(n/d)\lambda^{d}.
\]



\begin{thebibliography}{9}

\bibitem{posetslides}
A. Frieze,
\textit{Posets},
Combinatorics Lecture Slides, Carnegie Mellon University.
Available at: \url{https://www.math.cmu.edu/~af1p/Teaching/Combinatorics/Slides/Posets.pdf}

\bibitem{murty_combinatorics}
S. M. Cioab\u{a} and M. Ram Murty,
\textit{A First Course in Graph Theory and Combinatorics},
Springer / Hindustan Book Agency, 2009.

\end{thebibliography}

\end{document}
